% 由 generate-clustering-data.py 生成的真实 OPTICS 排序教学例。
% m=3（含自身），eps_max=1.1；无穷值用顶部三角形显示。
\begin{tikzpicture}
 \begin{axis}[
   width=\linewidth,height=27mm,at={(0,48mm)},anchor=south west,
   axis x line=bottom,axis y line=none,
   xmin=-0.4,xmax=12.5,ymin=-0.5,ymax=0.5,
   xtick={0,4,8,12},ytick=\empty,
   xlabel={原始特征值 $x$},font=\small,
   axis line style={BookMuted},tick label style={BookMuted}]
  \addplot[only marks,mark=*,mark size=1.8pt,BookInk]
   table[x=x,y expr=0] {figures/clustering-optics.dat};
 \end{axis}
 \begin{axis}[
   width=\linewidth,height=47mm,at={(0,0)},anchor=south west,
   xmin=0.4,xmax=23.6,ymin=0,ymax=1.15,
   axis lines=left,xlabel={OPTICS 访问次序},ylabel={距离},
   xtick={1,8,14,23},ytick={0,0.4,0.8},
   font=\small,axis line style={BookMuted},
   tick label style={BookMuted},unbounded coords=jump,
   legend style={draw=none,font=\footnotesize,at={(0.5,-0.42)},
     anchor=north,legend columns=3}]
  \addplot[ycomb,BookTeal,mark=*,mark size=1.3pt,line width=0.8pt]
   table[x=order,y=reach] {figures/clustering-optics.dat};
  \addlegendentry{可达距离}
  \addplot[BookBlue,dashed,line width=0.8pt]
   table[x=order,y=core] {figures/clustering-optics.dat};
  \addlegendentry{核心距离}
  \addplot[only marks,mark=triangle*,mark size=2pt,BookInk]
   table[x=order,y=restart] {figures/clustering-optics.dat};
  \addlegendentry{$+\infty$}
  \draw[BookAmber,dash dot] (axis cs:0.5,0.35) -- (axis cs:23.5,0.35);
  \node[BookAmber,anchor=south] at (axis cs:18,0.35) {$\varepsilon'=0.35$};
 \end{axis}
\end{tikzpicture}
